\chapter{Folic Acid (Folate)}
\section*{What Is This Test?}
Serum folate and red blood cell (RBC) folate measure the concentration of folate (vitamin B9), a water-soluble B vitamin essential for DNA synthesis, cell division, amino acid metabolism (homocysteine remethylation to methionine), and neural tube closure during embryogenesis (periconceptional folic acid supplementation reduces neural tube defects by 50--70\%). Folate, in the form of 5-methyltetrahydrofolate (5-MTHF), serves as the methyl donor for the remethylation of homocysteine to methionine (catalyzed by methionine synthase with vitamin B12 as a cofactor) \cite{bailey1999folate}. Folate deficiency impairs DNA synthesis in rapidly dividing cells, particularly hematopoietic cells in the bone marrow, producing megaloblastic anemia (macrocytic anemia with MCV >100 fL, hypersegmented neutrophils [$>$5 nuclear lobes in $>$5\% of neutrophils], and pancytopenia in severe deficiency). Serum folate reflects short-term folate intake (days to weeks) and is sensitive to recent dietary folate intake. RBC folate reflects total body folate stores over the preceding 2--3 months (the lifespan of erythrocytes) and is less affected by recent dietary intake. Causes of folate deficiency: inadequate dietary intake (most common --- alcoholics, elderly, poverty, poor diet), increased requirements (pregnancy, lactation, chronic hemolytic anemia, chronic inflammatory states), malabsorption (celiac disease, tropical sprue, extensive small bowel resection, bariatric surgery), and medications (methotrexate --- DHFR inhibitor; phenytoin, carbamazepine, valproic acid --- interfere with folate absorption/metabolism; trimethoprim, pyrimethamine --- DHFR inhibitors; oral contraceptives). Neural tube defects (NTDs --- anencephaly, spina bifida, encephalocele) are associated with maternal folate deficiency during the first 28 days of embryogenesis. Universal folic acid fortification of enriched cereal grains in the United States (mandated by the FDA in 1998) has reduced the incidence of NTDs by 25--50\%.
\section*{Alternative Names}
Folate; Folic Acid; Vitamin B9; Serum Folate; RBC Folate; Red Blood Cell Folate
\section*{Principle}
Serum and RBC folate are measured by competitive chemiluminescent immunoassay (CLIA) on automated immunoassay analyzers, using folate-binding protein (beta-lactoglobulin from cow's milk) as the specific binding agent.
\section*{History}
The role of folate in human disease was first recognized by Lucy Wills in 1931, who demonstrated that the macrocytic anemia of pregnancy in Bombay textile workers was cured by a yeast extract (Marmite), a factor she called the ``Wills factor'' \cite{wills1931treatment}. Folate itself was isolated from spinach leaves in 1941 by Mitchell and colleagues, who coined the name ``folic acid'' from the Latin folium (leaf). Victor Herbert's deliberate self-experiment on a folate-free diet in 1961 established the sequential hematologic and biochemical changes of folate depletion \cite{herbert1962experimental}. The landmark MRC Vitamin Study, published in 1991, proved that periconceptional folic acid supplementation reduces neural tube defects \cite{mrc1991prevention}, and US FDA-mandated fortification of enriched cereal grains followed in 1998. The folate--vitamin B12--homocysteine axis was clarified over the 1960s--1980s, linking both vitamins to methionine synthase and megaloblastic anemia.
\section*{Why This Test Is Ordered}
\begin{itemize}
  \item Evaluation of macrocytic anemia (MCV greater than 100 fL), especially with hypersegmented neutrophils or pancytopenia, to detect megaloblastic anemia
  \item Differentiation of folate deficiency from vitamin B12 deficiency, which produce indistinguishable megaloblastic changes on the CBC
  \item Assessment of folate status in patients with alcohol use disorder, malnutrition, restrictive diets, or dietary insufficiency \cite{debenois2008conclusions}
  \item Evaluation in pregnancy planning and in women with a prior neural tube defect-affected pregnancy, where periconceptional folate sufficiency is critical \cite{berry1999prevention}
  \item Investigation of malabsorptive states (celiac disease, tropical sprue, bariatric surgery, small bowel resection)
  \item Monitoring of patients on folate-antagonist drugs (methotrexate, trimethoprim, pyrimethamine, phenytoin, carbamazepine, valproic acid)
  \item Evaluation of unexplained hyperhomocysteinemia, in which folate deficiency is a common and treatable cause
\end{itemize}
\section*{Primary vs Secondary Test}
Folate testing is a secondary, confirmatory test in the evaluation of macrocytic or megaloblastic anemia, performed after the CBC suggests the diagnosis. It is ordered together with serum vitamin B12, and the distinction between the two deficiencies is finalized with serum methylmalonic acid and homocysteine levels.
\section*{How This Test Is Performed}
A blood sample is collected by standard venipuncture: serum folate is drawn into a plain (red-top) tube, and RBC folate into an EDTA (lavender-top) tube, from which the whole-blood folate is corrected by the hematocrit. Both are measured by competitive chemiluminescent immunoassay using folate-binding protein as the capture agent. Serum folate reflects intake over the preceding days to weeks, whereas RBC folate reflects average stores over the prior 2--3 months. Results are typically available within the same day.
\section*{What Preparation Is Needed}
An overnight fast (8--12 hours) is recommended before serum folate measurement because recent dietary folate intake can falsely raise the serum value; RBC folate is unaffected by recent intake. If feasible, blood should be drawn before folic acid supplements are started, since supplementation rapidly normalizes serum folate. Samples should be protected from light and processed promptly because folate is light-sensitive and degrades in prolonged storage.
\section*{Contraindications and Precautions}
There are no contraindications beyond those of routine venipuncture. Interpretive cautions: (1) Recent blood transfusion invalidates RBC folate because transfused donor cells carry donor folate levels. (2) Serum folate may be falsely normal shortly after a folate-rich meal and falsely low after recent alcohol intake; a low serum folate with borderline RBC folate warrants confirmation. (3) Hemolyzed specimens spuriously elevate measured folate. (4) Isolated low serum folate without anemia may reflect recent dietary variation rather than true tissue deficiency, so RBC folate or a repeat test is preferred for diagnosis. (5) Folate supplementation can partially correct the anemia of B12 deficiency while allowing neurologic damage to progress, so vitamin B12 status should always be evaluated before treating with folate alone.
\section*{Risks}
Minimal risk. Slight pain or bruising at the venipuncture site. Rarely: hematoma, infection, or vasovagal syncope.
\section*{What Happens After the Test}
Results are interpreted with serum vitamin B12, methylmalonic acid, homocysteine, and the CBC. Confirmed folate deficiency is treated with oral folic acid 1 mg daily, along with dietary counseling; in pregnancy, periconceptional supplementation of 400--800 mcg daily is standard. In megaloblastic anemia, a reticulocyte response is expected within 5--7 days of therapy, with hemoglobin correction over 4--8 weeks, so the CBC is usually repeated to confirm recovery.
\section*{Understanding Results}
\subsection*{Folate Deficiency (Serum <2--4 ng/mL or RBC <140 ng/mL)}
Megaloblastic anemia. Neural tube defects (periconceptional deficiency). Homocysteine elevation.
\subsection*{Folate Sufficiency (Serum >4 ng/mL, RBC >140 ng/mL)}
Adequate folate status.
\section*{Normal Results}
Serum folate: 2--20 ng/mL (>4 ng/mL = adequate, 2--4 ng/mL = borderline, <2 ng/mL = deficient). RBC folate: 140--960 ng/mL (>140 ng/mL = adequate). RBC folate is a better indicator of long-term folate status than serum folate. Fasting is recommended to avoid the confounding effect of recent food folate intake on serum folate.
\section*{Follow-up Tests}
Serum homocysteine (elevated in both folate and B12 deficiency), serum MMA (normal in folate deficiency, elevated in B12 deficiency) \cite{green2017vitamin}, serum vitamin B12, CBC (MCV, hemoglobin), anti-TTG IgA (celiac disease).
\section*{References}
\bibliographystyle{unsrtnat}
\bibliography{references}
\clearpage
